\documentclass[11pt]{article}

%% ============================================================
%% Comparison memo: the two existing CBS election-night models
%%   NYC 2025 general  (scripts/model/modeling-cory/nyc)
%%   TX 2026 primary   (scripts/model/modeling-cory/tx-prim-26)
%% Built from ~/Work/papers/memo-template.
%% Compile: pdflatex nyc_tx_comparison.tex  (twice, for refs)
%% ============================================================

\usepackage{amsmath}
\usepackage{mypaperfonts}
\usepackage{microtype}
\usepackage[margin=1in]{geometry}

\usepackage{graphicx}
\usepackage{booktabs}
\usepackage{array}
\usepackage{xcolor}
\definecolor{memolink}{RGB}{0,60,120}
\usepackage[hyphens]{url}
\usepackage[colorlinks=true, linkcolor=memolink, citecolor=memolink,
            urlcolor=memolink]{hyperref}
\usepackage{enumitem}
\setlist{noitemsep, topsep=2pt}

\usepackage{csquotes}

\usepackage[most]{tcolorbox}
\definecolor{memoalert}{RGB}{155,28,49}
\newtcolorbox{memobox}[2][]{%
  enhanced, breakable,
  colback=memoalert!4!white, colframe=memoalert!85!black,
  colbacktitle=memoalert!85!black, coltitle=white,
  fonttitle=\bfseries\small, fontupper=\small,
  title={#2},
  boxrule=0.6pt, arc=1.5pt, titlerule=0pt,
  left=7pt, right=7pt, top=6pt, bottom=6pt,
  before skip=10pt, after skip=10pt,
  #1}

\newcommand{\memoheader}[2]{%
  \noindent{\bfseries #1 \textperiodcentered\ Young-Kyung Kim \textperiodcentered\ #2}%
  \par\vspace{2pt}\hrule\vspace{1.2\baselineskip}%
}

\usepackage{titlesec}
\titleformat{\section}{\normalsize\bfseries}{\thesection.}{0.5em}{}
\titlespacing{\section}{0pt}{1.1\baselineskip}{0.35\baselineskip}
\titleformat{\subsection}{\normalsize\itshape\bfseries}{\thesubsection.}{0.5em}{}
\titlespacing{\subsection}{0pt}{0.9\baselineskip}{0.25\baselineskip}

\newcommand{\logit}{\operatorname{logit}}
\newcommand{\Norm}{\mathcal{N}}
\newcolumntype{L}[1]{>{\raggedright\arraybackslash}p{#1}}

\begin{document}

\memoheader{Two election-night models compared: NYC 2025 general vs.\ TX 2026 primary}{October 18, 2026}

\noindent Both pipelines in \texttt{scripts/model/modeling-cory} forecast an outcome from
partial precinct returns, both were run live, and they share an author and a utility
package. Architecturally they have almost nothing else in common. This memo sets them
side by side --- outcome, likelihood, prior, estimation, uncertainty --- and records what
each one measurably did on the night. All diagnostics quoted here were computed from the
committed code and its logged output; sources are given in \S\ref{sec:prov}.

\section{What each system was}\label{sec:what}

\begin{table}[ht]
\centering
\small
\begin{tabular}{@{}L{3.0cm}L{5.4cm}L{5.4cm}@{}}
\toprule
 & \textbf{NYC} & \textbf{TX} \\
\midrule
Contest & 2025 NYC mayoral, \emph{general} & 2026 TX \emph{primaries} (Senate, House) \\
Election day & 4 November 2025 & 3 March 2026 \\
Unit & election district (precinct) & precinct $\times$ vote method \\
Units modelled & 4{,}061 & 6{,}079 (Senate GOP) \\
Candidates & 3 (Mamdani, Cuomo, Sliwa) & up to 8 per contest \\
Live runs logged & 24, 21:31--00:27 EST & 21, 01:38--21:14 \\
Estimator lives in & \texttt{enightmodels} (binary dep.) & the repo itself \\
\bottomrule
\end{tabular}
\end{table}

\noindent One asymmetry is worth noting up front: TX's model is fully readable in the
repository, whereas NYC's core is a compiled dependency
(\texttt{CoryMcCartan/enightmodels}, pinned at \texttt{f3edb6e}) whose source tree is not
checked in. The only copy on disk is the installed build, readable via \texttt{lazyLoad}
on its database.

\section{Outcome and likelihood}\label{sec:like}

This is the deepest divergence, and most of the others follow from it.

\paragraph{NYC: Gaussian, turnout and share modelled separately.} Two kinds of equation,
sharing one control formula:
\begin{equation}
  \log(\tau_i + 0.1) = x_i'\beta^{\tau} + \varepsilon^{\tau}_i,
  \qquad
  s_{ij} = x_i'\beta_j + \varepsilon_{ij}, \quad j = 1,\dots,3 .
  \label{eq:nyc}
\end{equation}
The share equations are linear in the \emph{level} of the share --- no logit, no
multinomial --- and are fit independently. Predictions are clamped to $[0,1]$, observed
values substituted for reported precincts, and then normalised so the three shares sum to
one. That division is the only thing coupling the candidates.

\paragraph{TX: Poisson on counts, turnout implicit.} One model per candidate per method:
\begin{equation}
  \log \mathbb{E}[v_{icm}] = \beta_0 + \beta_1 \log(\widehat{\textrm{est}}_{icm} + 0.1)
  + f(\textrm{covariates}) + u_{\textrm{stratum}} + u_{\textrm{geostratum}} .
  \label{eq:tx}
\end{equation}
Because counts are modelled, turnout is implied ($\tau_{im} = \sum_j v_{ijm}$) rather than
estimated, and shares are formed afterwards as $(0.1 + \hat v)/(0.1 + \sum_j \hat v)$.
Note $\log(\widehat{\textrm{est}})$ enters as a free-coefficient covariate, not a fixed
offset, so the relationship to expected turnout is a fitted power law.

\begin{memobox}{On this axis, TX is structurally closer to right}
Modelling counts per candidate \emph{per mode} is implicitly a joint turnout-and-share
model, which matters because a precinct's overall share is a turnout-weighted mixture over
modes, $s_{ij} = \sum_m s_{ijm}\tau_{im} / \sum_m \tau_{im}$. If mail splits one way and
election day another, the precinct's share is determined by the \emph{mix}. NYC can treat
turnout and share as separable only because it drops mode entirely --- it reads
\texttt{vote\_mode} and then discards it, absorbing the effect into covariates such as
\texttt{prim25\_cvr\_eday} and an early-vote threshold. For a general election whose
reporting order is mode-driven, that is the weaker choice.
\end{memobox}

\section{The prior}\label{sec:prior}

\paragraph{NYC: a strong, informative $g$-prior.} The distinctive feature of the whole
pipeline. The prior is placed on \emph{fitted values}, not coefficients, and implemented by
stacking all $N$ precincts as pseudo-observations with weight $w = n_0/N$ on top of the
reported rows. The implied prior is exactly Zellner's $g$-prior,
\begin{equation}
  \beta \sim \Norm\!\big((X'X)^{-1}X'\tilde y,\; \sigma^2 (w X'X)^{-1}\big),
\end{equation}
with the location being the OLS projection of the prior fitted values. (``$g$-prior'' names
Zellner's structure, in which the prior covariance is proportional to $(X'X)^{-1}$; it is
not shorthand for ``Gaussian prior,'' and a ridge prior $\Norm(0,\lambda^{-1}I)$ is
Gaussian but not a $g$-prior. The proportionality to $X'X$ is what makes $n_0$ readable as
an effective sample size --- and also why no choice of $n_0$ improves the conditioning of
the design.) Those fitted values
come from a vendor likely-voter file (per-precinct \texttt{pred.turnout},
\texttt{pred.mamdani}, \texttt{pred.cuomo}, \texttt{pred.sliwa}, as of 9 October 2025),
rescaled by a single logit shift so that weighted aggregates hit exogenous targets:
borough turnout totals, and citywide shares of $0.47/0.38/0.12$.

\paragraph{TX: deliberately uninformative about candidates.} There is no candidate-level
prior at all. Each candidate's baseline is expected party turnout divided by the number of
candidates,
\begin{equation}
  \widehat{\textrm{est}}_{icm} = \textrm{est}_{im} \cdot \frac{\textrm{tr}}{n_{\textrm{cand}}},
\end{equation}
i.e.\ a flat split. Every candidate-level signal has to come from reported precincts. For a
primary with no precinct-level polling this is the right call, and it is a genuine design
difference rather than an omission.

\paragraph{Consequence: the two models move for different reasons.} In NYC the prior's
relative weight decays automatically as $n_0/(n_0 + n_{\textrm{obs}})$, so the forecast
walks from the vendor prediction toward the data. In TX there is nothing to walk from:
early-night behaviour rests entirely on $\log(\widehat{\textrm{est}})$ and a turnout-ratio
adjustment.

\section{Estimation}\label{sec:est}

\begin{table}[ht]
\centering
\small
\begin{tabular}{@{}L{3.6cm}L{5.1cm}L{5.1cm}@{}}
\toprule
 & \textbf{NYC} & \textbf{TX} \\
\midrule
Inference & conjugate Bayes, closed form & penalised ML \\
Machinery & one weighted QR per equation & \texttt{lme4::glmer}, Laplace approx. \\
Fallback & --- & \texttt{glmnet} Poisson, fixed $\lambda = 10$ \\
Returns & full posterior & point estimate \\
Draws & 500 per equation & none \\
Regularisation & $n_0$ pseudo-rows (hand-set) & random intercepts; lasso at $\lambda=10$ \\
Iterates? & no & yes (IRLS / coordinate descent) \\
\bottomrule
\end{tabular}
\end{table}

\noindent The Poisson log-likelihood $\sum_i\{v_i x_i'\beta - e^{x_i'\beta}\}$ is not
quadratic in $\beta$, so no finite-dimensional conjugate family exists. (Gamma conjugacy
applies to a free rate $\lambda_i$; imposing $\lambda_i = e^{x_i'\beta}$ with $\beta$
shared across precincts destroys it.) TX must therefore optimise, which yields a point
estimate and no posterior --- and that single fact explains its uncertainty machinery,
\S\ref{sec:unc}.

\begin{memobox}{The \texttt{glmnet} fallback has an uncontrolled prior}
\texttt{glm\_minnorm} calls \texttt{glmnet(X, y, family, lambda = 10)} with no
\texttt{alpha}, so it is \emph{lasso}. On a representative Poisson fit (Cornyn counts,
4{,}221 precincts, the model's own demographic covariates) $\lambda = 10$ zeroes five of
eight coefficients --- including \texttt{dem24\_est}, the 2024 partisan baseline, and both
race covariates --- against $\lambda_{\max} = 73.5$ for that data. The structural problem
is not the magnitude but that $\lambda$ is a fixed absolute number while $\lambda_{\max}$
is data-dependent: effective prior strength then varies unpredictably across contests and
across the night. At fifty-state scale this is unmanageable.
\end{memobox}

\section{Uncertainty}\label{sec:unc}

\paragraph{NYC: posterior draws, two intervals.} Coefficients are drawn as
$\beta^{(d)} = \hat\beta + R^{-1}z^{(d)}$ with $z^{(d)} \sim \Norm(0, \sigma^{(d)2}I)$ and
$\sigma^{(d)}$ from a scaled inverse-$\chi^2$; predictive noise resamples
leverage-corrected residuals $e_i/(1-h_i)$ once more than 25 precincts report. Aggregation
is draw-wise, $S_j = \sum_i s_{ij}\tau_i / \sum_i \tau_i$, so turnout and share uncertainty
propagate jointly. Both a credible interval and an anytime-valid confidence sequence are
reported at level $0.9999$.

\paragraph{TX: conformal residual quantiles.} With no posterior, intervals are built from
the cross-precinct spread of within-precinct share errors,
\begin{equation}
  r_{ic} = \frac{v_{ic}}{\sum_{c'} v_{ic'}} - \frac{\hat\mu_{ic}}{\sum_{c'}\hat\mu_{ic'}},
  \label{eq:conf}
\end{equation}
whose $\alpha/2$ and $1-\alpha/2$ quantiles across precincts are then added to
\emph{every} unreported precinct's share and summed.

\begin{memobox}{The TX interval assumes perfectly correlated precinct errors}
Adding the same quantile to all $\sim$6{,}000 precincts and summing is algebraically the
assumption that precinct-level errors move together as one common shock. A measure of
\emph{cross-sectional dispersion among} precincts is thereby reported as \emph{uncertainty
about their sum}. Under independence the aggregate error would scale as $\sigma_p/\sqrt{N}$,
roughly $77\times$ narrower; the truth lies between, but nowhere near $1\times$. Two
further problems in the same block: the residuals are \emph{in-sample} (predicted on the
very precincts used to fit), so there is no calibration split and no coverage guarantee ---
hence the \texttt{approx\_} label --- and \texttt{conf\_level} defaults to $0.95$ while the
output field is hardcoded \texttt{"approx\_conform\_99"}, so consumers are told 99\% while
95\% quantiles are computed.
\end{memobox}

\noindent The consequence is visible in the logs. For the Senate GOP contest the point
estimate was stable all night (Cornyn $40.6$--$43.9\%$), but the interval \emph{widened}
from $46.4$ to a peak of $63.5$ percentage points as counting progressed, because $q_{\min},q_{\max}$
are re-estimated each run from an increasingly heterogeneous conformity set and that growth
swamps the narrowing from substituted precincts. At the final logged run Cornyn's interval
was $[25.0,\,73.2]$ around a point estimate of $43.1\%$ --- fully overlapping Paxton's
$[21.3,\,69.8]$.

By contrast NYC's widths behave as theory predicts: the confidence sequence runs
$1.33$--$1.45\times$ the credible interval at the final snapshot, the premium one pays for
validity under roughly a hundred looks.

\section{Ensembles and specifications}\label{sec:ens}

\begin{table}[ht]
\centering
\small
\begin{tabular}{@{}L{3.0cm}L{5.4cm}L{5.4cm}@{}}
\toprule
 & \textbf{NYC} & \textbf{TX} \\
\midrule
Count & 4 specs, all reported & 7 specs in 4 fallback chains, + county legacy \\
Selection & none; consumers choose & \texttt{map\_until} takes the first that does not error \\
Complexity & 2 to 21 covariates & intercept-only to a 43-term spline \\
Prior strength & $n_0$ from $1.0$ to $15.5$, by eye & n/a \\
Nonlinearity & none, all linear & $\textrm{dem24\_est}^2$, thin-plate spline \\
Spatial & 64-term Fourier basis, \emph{disabled} & \texttt{geostratum} effects, active \\
Partial pooling & none & \texttt{stratum\_demo} $\times$ \texttt{geostratum} random effects \\
\bottomrule
\end{tabular}
\end{table}

\noindent NYC leans on many hand-picked covariates plus a strong prior; TX on few
covariates plus strata and splines with no candidate prior. For hundreds of races the TX
approach is the one that generalises --- one cannot hand-curate 21 covariates per contest
--- and its $5 \times 6$ strata cross is already the partial-pooling skeleton that sparse
legislative districts would need.

\section{Two measured problems specific to each}\label{sec:probs}

\paragraph{NYC: the prior is lossy in one spec, and silently so.} Because the $g$-prior
location is the projection of $\tilde y$ onto $\operatorname{col}(X)$, only the component
the covariates can express survives. Measuring the centred $R^2$ of each vendor prediction
on each spec's own column span:

\begin{table}[ht]
\centering
\small
\begin{tabular}{@{}lccc@{}}
\toprule
spec & Mamdani & Cuomo & Sliwa \\
\midrule
\texttt{simple\_pres} & $0.544$ & $\mathbf{0.073}$ & $0.711$ \\
\texttt{simple\_pred}, \texttt{simple\_mult} & $1.000$ & $1.000$ & $1.000$ \\
\bottomrule
\end{tabular}
\end{table}

\noindent Turnout is always safe, since $\log(\textrm{pred\_turn})$ is a covariate
everywhere. But \texttt{simple\_pres} does not contain the vendor share predictions, so its
prior on Cuomo's share is effectively a constant: $93\%$ of the vendor's precinct-level
Cuomo signal is annihilated by the projection. The four specs are therefore not a
$n_0$ ladder over a common prior --- they encode materially \emph{different} priors, and
since the prior predictive feeds the confidence sequence, that propagates into the
reported bands.

\paragraph{TX: two mechanisms never engaged.} The turnout-ratio adjustment filters on
\texttt{stage >= 6}, but the precinct feed emits only stages $0$, $2$ and $5$ (counts
$20/1511/2776$). The lookup table is therefore always empty and every run silently used the
hardcoded \texttt{turn\_ratio = 0.95}. Separately, \texttt{pct\_rep} in the logs is
non-monotonic --- $0.997 \to 0.269 \to 0.364$ --- so the feed regressed and the model's
reported base collapsed mid-count. Both are plumbing faults rather than modelling errors,
and both are invisible without reading the logs.

\section{Where each one's design should carry into 2026}\label{sec:impl}

\begin{table}[ht]
\centering
\small
\begin{tabular}{@{}L{5.6cm}L{3.3cm}L{3.3cm}@{}}
\toprule
Component & Take from & Rationale \\
\midrule
$g$-prior on fitted values, $n_0$ in precinct-equivalents & NYC & interpretable, decays correctly, exact conjugacy \\
Anytime-valid confidence sequence & NYC & the only honest interval under continuous monitoring \\
Draw-wise turnout-weighted aggregation & NYC & propagates both uncertainties jointly \\
Two-gate reporting logic (\texttt{obs} / \texttt{obs\_turn}) & NYC & separates ``votes in'' from ``count complete'' \\
Vote mode as a dimension of the unit & TX & the composition mechanism is structural \\
Strata + partial pooling & TX & required once districts are small \\
Small, fixed covariate set & TX & the only thing that scales to hundreds of races \\
\midrule
Post-hoc normalisation of independent share fits & \emph{neither} & deflates predictive variance; induces rather than estimates correlation \\
Conformal residual quantiles & \emph{neither} & assumes perfect cross-precinct correlation \\
Hardcoded $\lambda$ & \emph{neither} & effective strength varies with data scale \\
\bottomrule
\end{tabular}
\end{table}

\noindent The trial model in \texttt{tx\_gov\_trial.R} (documented separately in
\texttt{tx\_gov\_model.tex}) is an attempt at the top half of that table: a Gaussian,
conjugate, $g$-prior system in which vote mode is a dimension of the unit and the
composition is factored rather than normalised after the fact.

\section{Provenance of the diagnostics}\label{sec:prov}

\begin{itemize}
  \item NYC run log: \texttt{nyc/reports/data-out/nyc\_small\_history.json}, 24 records
        2025-11-05 02:31--05:27 UTC plus one at 19:48.
  \item NYC final snapshot and interval widths: \texttt{nyc\_current.json}.
  \item NYC projection losses: centred $R^2$ of each \texttt{pred\_*} column on each
        spec's design, computed on \texttt{nyc/data/nyc/nyc\_pre.parquet} (4{,}061 rows).
  \item NYC estimator internals: \texttt{forecast\_prec}, \texttt{lm\_bayes},
        \texttt{sim\_fitted}, \texttt{cred\_int}, \texttt{conf\_seq}, \texttt{fit\_t},
        read from the installed \texttt{enightmodels} lazy-load database.
  \item TX run log and interval widths: \texttt{tx-prim-26/data-out/log\_preds.ndjson},
        21 records 2026-03-04 01:38--21:14. Confirmed real data (\texttt{src = 0}): the
        logged times match none of the 31 hardcoded \texttt{fake\_times} on 2026-03-03.
  \item TX stage distribution: \texttt{tx-prim-26/data-out/precincts.csv}.
  \item \texttt{glmnet} shrinkage: Poisson fit on Cornyn counts from the same file.
\end{itemize}

\noindent One caveat on the NYC interval diagnostics: \texttt{fit\_t} computes
\texttt{MASS::fitdistr(x, "t")} and then discards the result, always returning
$(\textrm{mean}, \textrm{sd}, \textrm{df} = \infty)$. Both NYC intervals are therefore
Gaussian approximations rather than the $t$-fits the function name implies, and
\texttt{cred\_int} measures roughly $283\times$ slower than a plain quantile call entirely
on discarded computation. \texttt{cred\_int} also has no $[0,1]$ clamp, so a small
candidate's lower bound can print negative, though it did not at the final snapshot.

\end{document}
